% SPDX-License-Identifier: Apache-2.0
% covers: Dtm
\datasheet{Dtm}{The RISC-V debug transport over JTAG, behind the FPGA's BSCANE2, so that a stock OpenOCD, and gdb through it, reaches the core's debug module from the board's own cable.}

\dsfacts{\code{lib/parts/src/dtm.rs}}{\code{txhdl\_parts::dtm::Dtm}}%
{\code{//docs:examples}, \code{//docs:vreteno}}

\subsection*{Function}

The FPGA's JTAG port belongs to the device, so a design reaches it
through \code{BSCANE2}, one of the user scan chains of the device's
TAP. OpenOCD speaks to a transport there with its BSCAN tunnel,
\code{riscv use\_bscan\_tunnel 5}: it sets the device's instruction
register to \code{USER4}, and encapsulates each scan of the transport's own
TAP inside one data scan of that chain, framed as
OpenOCD 0.12.0's \code{riscv.c} frames it:

\begin{center}\footnotesize
\begin{tabular}{@{}lp{0.62\columnwidth}@{}}
\toprule
Scan & First bit shifted to last \\
\midrule
instruction & \code{0}, the width 5 in 7 bits, the register in 5,
\code{000} \\
data & \code{1}, the width $W$ in 7 bits, the data in $W + 1$,
\code{000} \\
\bottomrule
\end{tabular}
\end{center}

The payload is one bit longer than the register because what comes
back is a clock late: the device's TAP presents the bit that left on
the edge after. The unit shifts the first $W$ bits into the register
and lets the last go by.

The unit runs on \code{Tck}, the cable's clock, which \code{BSCANE2}
hands the fabric. Its request and answer cross to the system's clock
beside it, so it has one clock.

\subsection*{Registers}

The specification's, at its numbers.

\begin{center}\footnotesize
\begin{tabular}{@{}lllp{0.4\columnwidth}@{}}
\toprule
IR & Register & Width & What it is \\
\midrule
\code{0x01} & \code{idcode} & 32 & \code{0x00154001}, the transport's
own; selected after reset \\
\code{0x10} & \code{dtmcs} & 32 & version 1, \code{abits} 7, an idle
hint of one, \code{dmistat}; \code{dmireset} at bit 16 and
\code{dmihardreset} at bit 17 \\
\code{0x11} & \code{dmi} & 41 & address 7, data 32, op 2 \\
other & bypass & 1 & zero \\
\bottomrule
\end{tabular}
\end{center}

\dstables{Dtm}

\subsection*{Behaviour}

\begin{itemize}
\item A \code{dmi} scan that asks for a read (op 1) or a write (op 2)
sends its 41 bits on \code{req}, if nothing is in flight and no error
stands. The answer comes back on \code{ans}: the word above a two-bit
status, 0 for done and 2 for failed.
\item A scan that asks while an access is in flight, or that captures
\code{dmi} meanwhile, sets \code{dmistat} to 3, busy. OpenOCD answers
busy by idling longer before its next scan.
\item An answer of failed sets \code{dmistat} to 2.
\item Either error sticks, and no access goes while it stands, until
a write of \code{dtmcs} sets \code{dmireset}, which clears it, or
\code{dmihardreset}, which clears it and drops the access in flight.
\item The \code{reset} pin, the device's test-logic reset, selects
\code{idcode} again.
\item A \code{dmi} capture shows the address and the word of the last
access, with busy as its op while an access is in flight and
\code{dmistat} otherwise.
\end{itemize}

\subsection*{Verification}

\begin{itemize}
\item \code{//lib/parts:parts\_test} drives the device's side of
\code{BSCANE2} bit by bit: the \code{idcode} after reset,
\code{dtmcs}, bypass for an unknown instruction, a \code{dmi} write
then read, busy while an access is in flight until \code{dmireset},
and a failed access sticking.
\item \code{//lib/examples:ex\_dtm}, and the build simulates the
netlist against its run under nvc and Verilator
(\code{//docs:dtm\_sim\_dtm\_tb\_test}, \code{//docs:dtm\_vsim\_test}).
\item \code{//cpu/vreteno:openocd\_test} puts stock OpenOCD 0.12.0 on
the whole simulated board behind a model of the Artix-7's TAP: it
examines the hart, halts it, reads \code{pc} and memory, steps,
writes and resumes; the same \code{dtmcs} scans bit by bit as
OpenOCD's driver makes them; and gdb loading a program and stopping
at a breakpoint.
\end{itemize}

\subsection*{Used by}

The Vreteno board, on \code{USER4}, in every top. Its requests reach
the debug module through \code{DtmBridge}.

\subsection*{Limits}

One access in flight at a time, which is all the specification's
\code{dmi} has. Only the nested-TAP framing OpenOCD 0.12.0 uses by default is understood.
